\section{Diseño e implementación del sistema electrónico}

\subsection{Arquitectura general}

La plataforma separa tres funciones que la sección 8 exige mantener independientes: la adquisición, que debe ser determinista y local; el transporte, que puede ser diferido y tolerante a fallos; y el análisis, que debe ser reproducible a partir de los datos almacenados. La fuente instrumentada y cada nodo receptor se comportan como esclavos de una operación coordinada por un maestro. Cada nodo GEO convierte el movimiento del terreno en tensión, acondiciona y digitaliza la señal, la conserva localmente y la entrega cuando el enlace lo permite. Una interfaz de campo configura los nodos y gobierna las capturas sin requerir conectividad externa, mientras que un servidor centraliza registros, geometría y metadatos y ejecuta el flujo de procesamiento.

Debe señalarse desde ya que la arquitectura es multicanal por diseño, pero que la validación de campo disponible no lo es físicamente: cada registro se obtuvo con un nodo HAMMER y un único nodo GEO. La sección 10 explica esa restricción y delimita qué puede y qué no puede concluirse de ella.

\begin{figure*}[!tb]
\centering
\begin{tikzpicture}[
  caja/.style={draw,rounded corners=2pt,align=center,inner sep=3pt,
               minimum height=12mm,text width=21mm,font=\scriptsize},
  mues/.style={caja,draw=red!55!black,fill=red!4},
  det/.style={caja,draw=green!30!black,fill=green!4},
  com/.style={caja,draw=violet!65,fill=violet!5},
  sop/.style={caja,draw=blue!60,fill=blue!4},
  fl/.style={-{Latex[length=1.5mm]},semithick,gray!70!black},
  et/.style={font=\tiny,color=gray!60!black}]

  \node[mues] (sm) {\textbf{SM-24}\\[1pt]10\,Hz $\cdot$ $\zeta_0=0{,}25$};
  \node[mues,right=3mm of sm]  (afe) {\textbf{DA-AFE}\\[1pt]PGA\,$\to$\,BP\,$\to$\,SUM\,$\to$\,LP\\4 referencias ajustables};
  \node[det,right=3mm of afe]  (adc) {\textbf{ADC + DFB}\\[1pt]18 bits $\cdot$ 2604\,SPS\\FIR / estimador DC};
  \node[det,right=3mm of adc]  (udb) {\textbf{UDB + DMA}\\[1pt]superMaquina\\ruta determinista};
  \node[com,right=3mm of udb]  (esp) {\textbf{ESP32}\\[1pt]UART $\downarrow$ $\cdot$ I\textsuperscript{2}C $\uparrow$\\flanco local};
  \node[com,right=3mm of esp]  (mae) {\textbf{Maestro}\\[1pt]ESP-NOW $\cdot$ AP\\SPA + WebSocket};

  \draw[fl] (sm)--(afe); \draw[fl] (afe)--(adc); \draw[fl] (adc)--(udb);
  \draw[fl] (udb)--(esp); \draw[fl] (esp)--(mae);

  \node[sop,below=13mm of afe] (cal) {\textbf{Calibración \emph{foreground}}\\[1pt]AMUX + ADC + PI $\to$ códigos DAC\\se desactiva al adquirir};
  \node[sop,right=3mm of cal]  (mem) {\textbf{Memoria}\\[1pt]RAM ESP $\cdot$ RAM PSoC $\cdot$ SD\\crudo preservado};
  \node[sop,right=3mm of mem]  (rev) {\textbf{Revisión}\\[1pt]FastAPI $\cdot$ alineación\\control de calidad};
  \node[sop,right=3mm of rev]  (msw) {\textbf{MASW}\\[1pt]dispersión $\cdot$ \emph{backends}\\resultado: modo 0};

  \draw[fl] (cal) -- node[et,right=0.5mm]{antes de medir} (afe);
  \draw[fl] (udb.south) -- node[et,above right=-1mm and 0mm]{lotes} (mem.north);
  \draw[fl] (mem)--(rev); \draw[fl] (rev)--(msw);
\end{tikzpicture}
\caption{Arquitectura vigente del sistema y frontera entre el camino determinista de
muestra, en la fila superior, y las funciones de calibración, almacenamiento y análisis
que quedan fuera del instante crítico.}
\end{figure*}

\subsection{Transductor y el conflicto de banda}

\begin{figure*}[!tb]
\centering
\includegraphics[width=\textwidth,keepaspectratio]{respuesta_sm24_modelo.png}
\caption{Respuesta nominal del SM-24 en circuito abierto, con las formas referidas a velocidad y a aceleración en dos paneles de escala vertical propia. Bajo $f_n$ caen a razón de 40 y 20 dB por década respectivamente; son dos referencias de la misma tensión de salida, no dos modos de funcionamiento.}
\end{figure*}

El elemento sensor es un geófono vertical pasivo SM-24, de frecuencia natural nominal 10 Hz, sensibilidad 28,8 V·s/m y amortiguamiento dependiente de la carga \citep{SM24}. Su elección respondió a la disponibilidad institucional y no a una optimización de la profundidad objetivo, y es importante presentarla así porque de ella se derivan buena parte de las limitaciones del trabajo.

El conflicto con el requerimiento de banda de la sección 8 es cuantitativo. Un geófono se comporta como un sistema de segundo orden pasa-altos respecto de la velocidad del suelo, de modo que su respuesta cae aproximadamente 27,6 dB a 2 Hz y 40 dB a 1 Hz respecto de la zona plana situada por encima de su frecuencia natural. Como la frecuencia mínima necesaria para alcanzar 10 m es del orden de 7,5 Hz, el sensor opera justo en el codo de su propia respuesta; y como alcanzar 50 m requeriría trabajar en torno a 1,5 Hz, el objetivo formal cae de lleno en la región donde el transductor atenúa la señal cuarenta decibelios antes de que ninguna etapa electrónica pueda intervenir.

\subsection{Frente analógico y calibración}

\begin{figure*}[!tb]
\centering
\includegraphics[width=\textwidth,keepaspectratio]{afe_entrada.png}\\[3mm]
\includegraphics[width=\textwidth,keepaspectratio]{afe_compensa.png}
\caption{Cadena analógica implementada, en dos mitades del mismo esquemático.
Arriba: entrada
instrumental diferencial y etapa de ganancia programable.
Abajo: rama pasabanda de compensación en torno a \texttt{OPAbp}, con el
potenciómetro de 2\,k$\Omega$ que ajusta la cancelación, el sumador \texttt{OPAsum}
de 27\,k$\Omega$ y el pasa-bajos de realimentación múltiple \texttt{OPAlp}
(30\,k$\Omega$ / 150\,k$\Omega$ / 12\,k$\Omega$, 47\,nF y 3,3\,nF) previo al convertidor.}
\end{figure*}

El acondicionamiento implementado en el PSoC 5LP encadena una entrada instrumental, una etapa de ganancia programable, una rama de compensación, un sumador y un filtrado pasa-bajos previo a la conversión. La etapa de compensación es la decisión de diseño central del frente analógico y merece desarrollarse con detalle.

El punto de partida es que una misma tensión de salida del geófono admite dos lecturas. Referida a la velocidad del suelo, la respuesta del sensor es un pasa-altos de segundo orden,

\begin{equation}
H_v(s)=\frac{V_o(s)}{\dot X_g(s)}=-G_0\,\frac{s^{2}}{s^{2}+2\zeta_0\omega_n s+\omega_n^{2}},
\end{equation}

donde $G_0$ es la sensibilidad, $\omega_n=2\pi f_n$ la pulsación natural y $\zeta_0$ el amortiguamiento nominal. Como la aceleración cumple $A_g(s)=s\dot X_g(s)$, la misma salida referida a aceleración resulta

\begin{equation}
H_a(s)=\frac{H_v(s)}{s}=-G_0\,\frac{s}{s^{2}+2\zeta_0\omega_n s+\omega_n^{2}},
\end{equation}

que es un pasabanda y no presenta meseta con el amortiguamiento nominal. Por debajo de $f_n$, $|H_v|$ cae a razón de 40 dB por década y $|H_a|$ a 20 dB por década. No se trata de dos modos de funcionamiento sino de dos referencias de la misma tensión, y la distinción importa porque la estrategia de compensación se formula sobre la referencia a aceleración.

La observación que aporta Ma y colaboradores es que no hace falta invertir la respuesta del sensor, operación mal condicionada porque amplificaría el ruido con la misma pendiente con que el geófono atenúa la señal. Basta con conservar $f_n$ y elevar el amortiguamiento efectivo hasta un valor $\zeta_1\gg\zeta_0$, con lo cual la respuesta referida a aceleración se aplana entre

\begin{equation}
f_1\simeq\frac{f_n}{2\zeta_1},
\qquad
f_2\simeq 2\zeta_1 f_n,
\end{equation}

de modo que $\zeta_1$ se elige para llevar el quiebre inferior $f_1$ hasta la banda de interés. El quiebre superior $f_2$ no constituye un objetivo, ya que la aplicación no requiere extender la banda alta; sólo se conserva la simetría logarítmica $f_1f_2=f_n^{2}$ en torno a $f_n$.

La transferencia que realiza esa transformación es el cociente entre el denominador nominal y el deseado,

\begin{equation}
H_{\mathrm{comp}}(s)=\frac{s^{2}+2\zeta_0\omega_n s+\omega_n^{2}}{s^{2}+2\zeta_1\omega_n s+\omega_n^{2}},
\end{equation}

y al separar el término constante aparece la forma que da nombre a la arquitectura y que se implementa directamente en el circuito,

\begin{equation}
H_{\mathrm{comp}}(s)=1-\frac{2(\zeta_1-\zeta_0)\,\omega_n s}{s^{2}+2\zeta_1\omega_n s+\omega_n^{2}}.
\end{equation}

La estructura $1-\mathrm{BP}$ no es entonces una elección heurística sino una consecuencia algebraica: el camino directo menos un pasabanda centrado en la resonancia del sensor. Su implementación se resolvió desacoplando los dos polos, uno en la entrada y otro en la realimentación de un inversor, de manera que cada uno dependa de un único par $RC$. Esa realización se comparó experimentalmente contra una Sallen-Key y contra una de realimentación múltiple. La Sallen-Key quedó descartada porque su realimentación positiva agrava la sensibilidad del $Q$ a las tolerancias y admite oscilación si los componentes se apartan del valor nominal, mientras que la desviación medida respecto de la respuesta objetivo resultó menor en la realimentación múltiple. Conviene precisar que esta última no se abandonó: es la topología con la que está implementado el filtro pasa-bajos antialias que precede al convertidor. Los valores se resolvieron por búsqueda sobre la malla de componentes comerciales incorporando el modelo no ideal del amplificador del PSoC, y un potenciómetro en serie con la rama pasabanda permite ajustar en cada unidad la proporción de la que depende la cancelación.

\begin{figure*}[!tb]
\centering
\includegraphics[width=\textwidth,keepaspectratio]{extension_banda.png}
\caption{Efecto de la compensación sobre la respuesta referida a aceleración. Al elevar $\zeta_1$ los polos complejos se separan en dos reales, y el inferior desplaza el quiebre $f_1$ hacia la banda de interés.}
\end{figure*}

El precio de la compensación es alto y conviene explicitarlo, porque condiciona todo el resto del frente analógico. Con los valores implementados se obtiene $f_0=10{,}20$ Hz y un amortiguamiento efectivo $\zeta_1\simeq 937$, que es el valor con el que se diseñó la rama pasabanda. En la zona central la respuesta compensada queda escalada por $1/(2\zeta_1)\simeq 5{,}3\times10^{-4}$, esto es, una atenuación cercana a 1875 veces que debe recuperarse íntegramente en ganancia. Esa ganancia se repartió deliberadamente entre etapas para no concentrar el margen en un solo bloque: la entrada instrumental aporta $\times 2$, la etapa programable entre $\times 1$ y $\times 50$, el sumador $\times 3{,}6$ sobre la rama de compensación y el pasa-bajos $\times 5$, con un producto máximo de $\times 1800$. El propósito de la compensación no es convertir al SM-24 en un sensor de banda ancha, objetivo inalcanzable por vía puramente electrónica, sino conservar una forma de respuesta predecible en la banda baja y aprovechar el margen del convertidor, aceptando a cambio un compromiso explícito con el ruido: toda la cadena de ganancia que restituye ese factor 1875 amplifica también el ruido propio del frente, y es esa tensión la que vuelve indispensables las mediciones de ruido referido a entrada que la sección 12 declara pendientes.

Extender la banda hacia baja frecuencia agrava el problema del offset y de la deriva, que de otro modo consumirían el rango del convertidor sin aportar señal. Para controlarlo, un multiplexor permite observar nodos internos de la cadena con el mismo convertidor, y una rutina de calibración en primer plano ajusta secuencialmente las referencias de la etapa de ganancia, de la rama de compensación, del sumador y del filtro. Durante la captura los lazos se desactivan y se aplican códigos fijos almacenados con verificación de integridad, de manera que ninguna corrección actúe sobre la señal mientras se registra. La revisión posterior del circuito, que sustituye las referencias por fuentes de corriente e incorpora una etapa de ganancia de salida, se encuentra en desarrollo y no participó de la campaña; los resultados de la sección 10 corresponden exclusivamente a la ruta anterior.

\begin{figure*}[!tb]
\centering
\includegraphics[width=0.62\textwidth,keepaspectratio]{flujo_calibracion_horizontal_usuario.png}
\caption{Flujo de la calibración \emph{foreground}: las etapas se ajustan en orden, cada una hasta entrar en su banda muerta, y la adquisición se restaura con códigos fijos.}
\end{figure*}

\subsection{Digitalización y temporización determinista}

\begin{figure*}[!tb]
\centering
\begin{minipage}[t]{0.53\textwidth}\centering\includegraphics[width=\textwidth,keepaspectratio]{dig_control.png}\end{minipage}\hfill
\begin{minipage}[t]{0.44\textwidth}\centering\includegraphics[width=\textwidth,keepaspectratio]{dig_dma.png}\end{minipage}
\caption{Camino digital implementado en el PSoC. Izquierda: bloque de control
\texttt{SignalController} descrito en lógica programable, con los registros de
configuración, el filtro FIR por hardware y el temporizador de vigilancia.
Derecha: registros de estado y los tres canales de acceso directo a memoria que
trasladan la señal cruda y la filtrada sin intervención del procesador.}
\end{figure*}

La conversión emplea el convertidor delta-sigma del PSoC configurado a una frecuencia de muestreo del orden de 2604 muestras por segundo, con resolución nominal de 18 bits y rangos bipolares seleccionables que se coordinan con la ganancia programable para aprovechar la escala sin saturar. Un filtro FIR implementado en hardware y varios canales de acceso directo a memoria trasladan la señal cruda y la filtrada sin intervención del procesador.

La elección de una arquitectura delta-sigma frente a una de aproximaciones sucesivas responde a la forma particular del problema. La banda de interés termina por debajo de unos pocos cientos de hertz, de modo que la velocidad de conversión no es una restricción, mientras que el rango dinámico sí lo es: el mismo registro debe contener el primer arribo cercano a la fuente y la llegada atenuada del offset más lejano. Un convertidor delta-sigma explota precisamente ese desbalance, ya que sobremuestrea muy por encima de la banda útil y desplaza el ruido de cuantificación hacia frecuencias altas, donde el diezmado lo elimina; obtiene así una resolución que una arquitectura de aproximaciones sucesivas sólo alcanzaría con un comparador y una referencia considerablemente más costosos. El sobremuestreo aporta una segunda ventaja relevante en campo: relaja la exigencia sobre el filtro antialias analógico, cuyo orden y precisión de componentes habrían tenido que aumentar para proteger una conversión directa a esa velocidad. A ello se suma que el convertidor está integrado en el mismo dispositivo que aloja el acondicionamiento y la lógica de control, lo que evita una interfaz externa y sus problemas asociados de acoplamiento y de temporización.

La temporización de la ventana de captura es el punto donde la sección 8 impone su condición más estricta, y por eso no se delegó al software. Una máquina de estados descrita en lógica programable coordina el disparo, la habilitación del convertidor, el conteo de muestras y el cierre de la captura, de modo que la duración y el instante de la ventana no dependen de la latencia variable del firmware. Esta decisión es la que hace defendible el requerimiento de sincronización, aunque su verificación experimental siga pendiente. Queda además una decisión de diseño abierta que incide directamente sobre esa métrica: la selección de la fuente de reloj que gobierna la ventana. El oscilador interno del dispositivo resulta suficiente para fijar la duración de la captura, pero su deriva y su incertidumbre de frecuencia se acumulan sobre registros largos y entre nodos distintos, de modo que evaluar un cristal externo o una referencia compartida es parte del trabajo necesario para acercarse a la tolerancia derivada en la sección 8.

Corresponde una advertencia explícita sobre la resolución. Los 18 bits son una cifra nominal del convertidor y no un número efectivo de bits medido: faltan ensayos de ruido referido a entrada, linealidad, saturación y repetibilidad por rango y por ganancia. Que la campaña haya registrado señales utilizables demuestra que la cadena funciona, no que su desempeño esté cuantificado.

\subsection{Sincronización, comunicaciones y trazabilidad}

La coordinación entre nodos y el transporte de datos se resuelven sobre módulos ESP32 mediante un protocolo de difusión directa entre pares, que evita depender de infraestructura de red en campo. Como una captura completa excede ampliamente la carga útil de un paquete y el enlace no garantiza entrega, el protocolo fragmenta cada registro, numera las secuencias, confirma la recepción y reintenta los fragmentos faltantes, mientras cada nodo conserva además una copia local. El motivo de exigir la entrega íntegra, y no una reconstrucción aproximada, es que la magnitud que el método estima es una fase: un tramo de muestras faltante no degrada gradualmente el resultado sino que invalida el registro completo para el cálculo de dispersión. De ahí que la política ante una pérdida sea reintentar o descartar, y que la copia local exista precisamente para que descartar no signifique perder la posición.

La separación entre la interfaz de campo, que opera los nodos, y el servidor, que valida, agrupa y preserva los registros junto con su geometría y sus metadatos, es la que permite reconstruir a posteriori cualquier resultado a partir de los datos almacenados.

\subsection{Fuente instrumentada}

La excitación se aplica con un martillo instrumentado PCB 086D20, cuyo sensor piezoeléctrico de tipo ICP entrega una medida de la fuerza aplicada y requiere alimentación por corriente constante con desacople de la polarización \citep{PCB086D20,PCBSignalConditioning}; el acondicionamiento correspondiente se integró en el mismo PSoC. El nodo HAMMER y el nodo GEO se adquieren de forma conjunta y con arranque coordinado, lo que permite referenciar cada respuesta al instante del impacto y, sobre todo, apilar varios golpes por posición.

La repetibilidad de la fuente adquiere una importancia particular en el protocolo empleado. Cuando el arreglo se construye desplazando un receptor, cualquier variación de amplitud, punto de golpe o acoplamiento entre posiciones sucesivas se incorpora al registro como si fuera variación espacial. Registrar entre veinte y treinta y siete golpes por posición y seleccionar el subconjunto más consistente reduce ese riesgo pero no lo elimina, y por esa razón la geometría aparece en la sección 11 como la sensibilidad dominante del resultado.


\begin{figure*}[!tb]
\centering
\includegraphics[width=\textwidth,keepaspectratio]{repetibilidad_fuente.png}
\caption{Repetibilidad del impacto medida sobre el canal del martillo. Los puntos claros son golpes individuales y las líneas la mediana por posición.}
\end{figure*}